\section{The sharply flag-transitive and nonspread branches}
We retain the projective flag-transitivity classification in source
Proposition 3.1, attributed to Kantor and Cameron--Kantor: besides the
classical group and the $A_7$ action on $\operatorname{PG}(3,2)$, a
flag-transitive projective plane has prime point number $m^2+m+1$ and
a sharply flag-transitive Frobenius group of order
$(m^2+m+1)(m+1)$. We also retain Feit's restrictions in Proposition 3.2:
the Desarguesian possibilities are $m=2,8$, while another example
requires $m\equiv0\pmod8$ and
$d^{m+1}\equiv1\pmod{m^2+m+1}$ for every divisor $d$ of $m$.
The reported computer extension of Feit's bound to $m\le10^{10}$ is
not a premise of the following proof and has no supplied source program.

\begin{proposition}[The exceptional projective branch]\label{projectivefrobenius}
A locally primitive sharply flag-transitive projective plane in this
retained branch is the Fano plane with group $7:3$. It is not locally
two-transitive.
\end{proposition}
\begin{proof}
The point stabilizer acts regularly on its $m+1$ incident lines.
A regular action is primitive only in prime degree, so $m+1$ is
prime. For a putative non-Desarguesian example, put $N=m^2+m+1$.
We have $m^3\equiv1\pmod N$, and the order of $m$ modulo the prime
$N$ is exactly three, since $1<m<N$. Feit's condition with $d=m$
therefore gives $3\mid m+1$, forcing $m=2$, a contradiction.
Among the Desarguesian possibilities, only $m=2$ has $m+1$ prime.
The group $7:3$ is regular on flags of the Fano plane; both local
three-point groups are cyclic regular. They are primitive but are
not two-transitive.
\end{proof}

The other exceptional projective action has $A_7$ on fifteen points
and thirty-five lines of $\operatorname{PG}(3,2)$. Retain its classical
identification and point stabilizer $\operatorname{PSL}_3(2)$ from
Cameron--Kantor/Taylor, as in source Proposition 3.3. The pencil is its
natural seven-point plane action, two-transitive by the elementary
basis argument. The stabilizer of a line induces $S_3$ on its three
points. With these exact retained local identifications, its claimed
local two-transitivity follows. This geometric action must not be
confused with the circle on the same fifteen points.

\begin{proposition}[One-dimensional nonspread actions, including the small boundary]\label{nonspread}
Retain Delandtsheer's translation/nonspread reduction quoted in source
Propositions 5.2 and 5.7. A flag-transitive one-dimensional nonspread
linear space is locally primitive precisely in the following two cases:
the three-point circle with $S_3=\operatorname{AGL}_1(3)$, and the
Fano plane with $7:3$. Only the first is locally two-transitive.
\end{proposition}
\begin{proof}
Write $v=p^d$ and let $k$ and $r=(v-1)/(k-1)$ be line size and pencil
size. The point stabilizer is a subgroup of $\Gamma L_1(p^d)$,
cyclic by cyclic. Its transitive action is primitive exactly when
$r$ is prime, by Lemma~\ref{prime} and the elementary fact that a
transitive prime-degree action is primitive. The retained nonspread
reduction gives $v\mid b$, where $b=vr/k$, and hence $k\mid r$.
Thus $k=r$ and $v=k(k-1)+1$.

There are now $b=v$ lines. Every other line meets a fixed line $L$:
its $k$ points lie on $k-1$ other lines each, giving $k(k-1)=v-1$
distinct lines, since two distinct points of $L$ cannot share another
line. If $k=2$ this is exactly the three-point circle. Flag transitivity
requires all six permutations, so the specified group is $S_3$; both
local actions have degree two.

If $k\ge3$, this is a projective plane of order $k-1$. The retained
projective classification and the preceding proposition leave exactly
the Fano plane and its group $7:3$. Conversely these two structures
have the asserted one-dimensional realizations and primitive local
actions. The cyclic local groups in the Fano case are not two-transitive.
\end{proof}

The source's Proposition 5.8 omits the $k=2$ case when passing from
$r=k$ to a nondegenerate projective plane. Its earlier Proposition 5.1
already lists this three-point circle. The footnote after Proposition 5.6
saying the spread restriction can be removed must therefore include this
extra three-point action. The final circle branches of Theorems 6.1
and 6.2 already include it as the natural degree-three $\operatorname{PSL}_2(2)$
action. No new family beyond that final union is being claimed.
